%! Two Categorical Variables — Unit 2, Lesson 2.1 — Student Packet Cover Sheet
\documentclass[12pt]{article}
\usepackage{apstats-article}
\usepackage{apstats-boxes}
\usepackage{ltablex}
\keepXColumns

\begin{document}

% Full-bleed navy banner. \coverbanner MEASURES the title block and sizes the
% band to it, so a lesson title long enough to wrap grows the band rather than
% running out the bottom of it.
\coverbanner{Unit 2}{Lesson 2.1 \quad Two Categorical Variables}

% NAMESTRIP: this is the ONLY component that carries the name row.
\namedateperiod
\vspace{0.18in}

\boxguard[14]
\begin{learningtargetbox}
\small
% GRADUAL RELEASE: the vocabulary is taught in the guided notes, so the targets name it
% plainly. The AP Learning Objectives (VAR-1.D, UNC-1.P, UNC-1.Q, UNC-1.R) stay in the plan.
\begin{enumerate}[leftmargin=1.5em, itemsep=5pt]
  \item \ldots{} ask whether two variables are related, and read a \textbf{two-way table} of
        two categorical variables.
  \item \ldots{} find \textbf{joint}, \textbf{marginal}, and \textbf{conditional relative
        frequencies}, choosing the right total to divide by.
  \item \ldots{} read a \textbf{segmented bar graph} and decide whether two variables show an
        \textbf{association}, justified with percents in context.
  \item \ldots{} explain why an association in a survey does not show that one variable causes
        the other.
\end{enumerate}
\end{learningtargetbox}

\vspace{0.14in}

\boxguard[14]
\begin{tocbox}
\small
\renewcommand{\arraystretch}{1.5}
\begin{tabularx}{\linewidth}{@{} c l X r @{}}
\rowcolor{sky}
\textbf{\#} & \textbf{Component} & \textbf{Description} & \textbf{Score} \\
\midrule
1 & Warm-Up & A relative frequency, two quiz sections, and a question about related
                 variables & \blank{1.2cm} \\
2 & Guided Notes \& Practice & Phones and sleep in a two-way table: joint, marginal, and
                 conditional relative frequencies, then texting drivers together & \blank{1.2cm} \\
% AP Practice is assigned but NOT scored: its score cell prints NA, never a \blank{}.
% Row 4 is the lesson's GRADED practice. Paper homework is the DEFAULT for this course; on the
% occasions DeltaMath is assigned instead, that replaces row 4 and its score cell is struck.
3 & AP Practice & A grocery survey: two multiple-choice items and one free-response set ---
                 practice, not a grade & \textbf{NA} \\
4 & Homework & A gym's membership records: read the table, then judge the association ---
                 \textbf{scored, due the first class after two study halls} & \blank{1.2cm} \\
\midrule
  & \multicolumn{2}{r}{\textbf{Total}} & \blank{1.2cm} \\
\end{tabularx}
\end{tocbox}

\vspace{0.14in}

\boxguard[8]
\begin{remindbox}
\small
Before you divide, ask \textbf{``of whom?''} The words after \emph{of those who} name the group,
and that group's total goes on the bottom. When two groups are different sizes, the bigger count
can hide the smaller rate --- \textbf{compare percents, not counts}. And a real association in a
survey still is not proof that one variable \emph{causes} the other.
\end{remindbox}

\end{document}
